\chapter{Harmonogram realizacji projektu}
\label{ch:harmonogram}

Prace nad projektem prowadzone by\l{}y w~trakcie semestru letniego roku
akademickiego 2025/2026. Tabela~\ref{tab:harmonogram} przedstawia podzia\l{}
zada\'{n} na etapy wraz z~przybli\.{z}onym czasem realizacji.

Rysunek~\ref{fig:gantt} przedstawia te same etapy w~formie diagramu Gantta.

\begin{figure}[H]
  \centering
  \includegraphics[width=\textwidth]{img/gantt}
  \caption{Diagram Gantta harmonogramu realizacji projektu}
  \label{fig:gantt}
\end{figure}

\begin{table}[H]
\centering
\caption{Harmonogram realizacji projektu}
\label{tab:harmonogram}
\begin{tabularx}{\linewidth}{|c|X|c|}
\hline
\textbf{Etap} & \textbf{Zakres prac} & \textbf{Czas (tygodnie)} \\
\hline
1 & Analiza wymaga\'{n}, dobor technologii, konfiguracja \'{s}rodowiska & 1 \\
\hline
2 & Projekt bazy danych, definicja modeli EF Core, pierwsze migracje & 1 \\
\hline
3 & Implementacja warstwy danych i~ViewModeli kursant\'{o}w, instruktor\'{o}w
    i~pojazd\'{o}w & 2 \\
\hline
4 & Implementacja kalendarza rezerwacji (ViewModel + widok XAML) & 2 \\
\hline
5 & Implementacja logowania i~kontroli dost\k{e}pu ról & 1 \\
\hline
6 & Testy funkcjonalne, naprawa b\l{}\k{e}d\'{o}w, refaktoryzacja & 1 \\
\hline
7 & Bezpiecze\'{n}stwo: bcrypt, walidacja PESEL, TOCTOU, indeksy unikalne & 1 \\
\hline
8 & Pisanie dokumentacji, ko\'{n}cowa weryfikacja & 1 \\
\hline
\multicolumn{2}{|r|}{\textbf{\l{}\k{a}cznie:}} & \textbf{10} \\
\hline
\end{tabularx}
\end{table}